\documentclass{article}
\begin{document}
\title{A Demonstration Manuscript with Planted Problems}
\maketitle

\section{Results}
The $N=18$ participants gave a mean usefulness rating of 5.19 on a 1--7 scale.
Our model reached an accuracy of 0.847 (Table~\ref{tab:main}).
The ensemble outperformed the strongest baseline by 1.4 points (95\% CI $-0.3$ to 3.1).
The difference was significant ($t = 1.20$, $p = 0.003$).
Calibration improved across horizons (Figure~\ref{fig:cal}d), with Brier scores listed in the released results.

\begin{table}[h]
\caption{Main results.}\label{tab:main}
\begin{tabular}{lcc}
Method & Accuracy & F1 \\
Our model & 0.851 & 0.802 \\
Baseline & 0.812 & 0.774 \\
\end{tabular}
\end{table}

\begin{figure}[h]
\caption{(a) Reliability. (b) Drift. (c) Horizon error.}\label{fig:cal}
\end{figure}
\end{document}
